\paragraph{A contraction gate involving only quadratic-many factors.}
The unresolved mass comparison can be reduced further to a
dimension-increment inequality. This reduction proves the error
from the two selected sibling branches; it does not assume or
prove the remaining uniform contraction inequality.

For $n\ge4$ write $M_n=2^{n-1}$, let
$b_n=3^{-1}\pmod {M_n}$, and use the positive multiset
$\mathcal S_n$ of (P-F56u). For every odd residue $a\bmod M_n$
define
\[
 P_n(a)=
 \prod_{b\in\mathcal S_n}
       |1-\exp(2\pi iab/M_n)|^{-1},\qquad
 Q_n=\sum_{\substack{a\bmod M_n,\ a\ {\rm odd}\\
                               a\ne b_n,-b_n}}
                              \frac{P_n(a)}{P_n(b_n)}.
 \tag{P-F56aa}
\]
Every factor is nonzero. Multiplication of the primitive-root
index by the odd number $g=\gcd(n,3)$ permutes the odd residues.
Thus (P-F56v) implies exactly $Q_n=E_n/|Z_n|$; the common
normalization $1/M_n$ cancels.

For $n\ge5$, form the multiset of residues modulo $M_n$
\[
 \begin{split}
 \mathcal T_n={}&\{1\}\\
 &\sqcup\{1+w_j+w_k:2\le j<k<n,\ (j,k)\ne(2,3)\}\\
 &\sqcup\{2+w_k:3\le k<n\}\sqcup\{6\},
 \qquad |\mathcal T_n|=\binom{n-1}2,
 \end{split}
 \tag{P-F56ab}
\]
where $w_i=(-2)^i$ for $i<n-1$ and $w_{n-1}=0$,
as before. Define the positive increment product
$J_n(a)=\prod_{b\in\mathcal T_n}
 |1-\exp(2\pi iab/M_n)|^{-1}$.
The exact recursion is
\[
             P_n(a)=J_n(a)P_{n-1}(-a\bmod M_{n-1}).
 \tag{P-F56ac}
\]

Here is a sufficient uniform contraction hypothesis, involving
only the $\binom{n-1}2=O(n^2)$ factors of each increment:
there exist a fixed $0<\rho<1$ and an integer $n_0\ge5$ such that,
for every $n\ge n_0$ and every odd $a\bmod M_n$ with
$a\bmod M_{n-1}\notin\{b_{n-1},-b_{n-1}\}$,
\[
 J_n(a)+J_n(a+M_{n-1})\le \rho J_n(b_n).
 \tag{P-F56ad}
\]
This inequality remains unproved. If it holds, then
$Q_n\longrightarrow0$, hence (P-F56n) holds for every
sufficiently large $n$ and gives the conditional exponential
index conclusion (P-F56p).
In this implication there is no unproved finite base
certificate: any finite starting value of $Q_n$ suffices.

The selected sibling error is bounded unconditionally:
\[
 \delta_n:=
 \frac{P_n(b_n+M_{n-1})}{P_n(b_n)}
 \le
 2^{-n^3/6+4n^2-107n/6+19},
 \qquad
 \delta_n\le2^{-n^3/12}\quad(n\ge48).
 \tag{P-F56ae}
\]
The exact implication of (P-F56ad) is
\[
                         Q_n\le\rho Q_{n-1}+2\delta_n.
 \tag{P-F56af}
\]

\emph{Proof.}
Remove the first variable from the chain weights.
The remaining weights $w_1,\ldots,w_{n-1}$ are $-2$
times the weights of the $(n-1)$-variable chain, in the
same order. In the multiset (P-F56u), all triples with
smallest index at least one occur. The shifted smaller
multiset removes from these exactly the triples
$\{1,2,k\}$ for $3\le k<n$ and $\{1,3,4\}$.
The additional triples with smallest index zero are
$\{0,j,k\}$ for $2\le j<k<n$, except $\{0,2,3\}$.
The additional singleton is $w_0=1$.
These extra weights are (P-F56ab): the first removed
family has weight $w_1+w_2+w_k=2+w_k$, and the last
triple has weight six modulo $M_n$.
For $n\ge6$ its integer weight is $-2-8+16=6$;
for $n=5$ it is $-2-8+0=-10\equiv6\pmod {16}$.
Thus the multiset decomposition is exact, including
multiplicities and the smallest allowed dimension.
Its cardinality is
$1+(\binom{n-2}2-1)+(n-3)+1=\binom{n-1}2$.
All its factors are nonzero because it is a submultiset
of $\mathcal S_n$. Substitution of the root
$\exp(-4\pi ia/M_n)$ in the smaller product proves
(P-F56ac).

Put $M=M_n$ and $\eta=\exp(2\pi ib_n/M)$.
The two roots indexed by $b_n$ and $b_n+M/2$
have the same smaller-dimensional factor in (P-F56ac).
In the increment, every even weight gives the same
factor at $\eta$ and $-\eta$.
The odd weights of $\mathcal T_n$ are exactly
\[
 \mathcal O_n=
 \{1\}\sqcup
 \{1+w_j+w_k:2\le j<k<n,\ (j,k)\ne(2,3)\}.
 \]
It follows that
\[
                 \delta_n
       =\prod_{b\in\mathcal O_n}
                     \left|\tan\frac{\pi b_n b}{M}\right|.
 \tag{P-F56ag}
\]
No tangent is singular, since $b_n$ and every $b$ here
are odd and $M\ge16$.

Among these odd weights, the bulk pairs
$2\le j<k\le n-2$, excluding $(2,3)$,
give
$B=\binom{n-3}2-1$ weights divisible by three.
For such a weight $b=1+w_j+w_k$, one has
$|b|<2^{k+1}$ and $|b|\le M$.
Since $3b_n\equiv1\pmod M$,
$\tan(\pi b_n b/M)=\tan(\pi b/(3M))$.
On $[-\pi/3,\pi/3]$, the bounds
$|\sin u|\le|u|$ and $\cos u\ge1/2$ give
\[
 \left|\tan\frac{\pi b_n b}{M}\right|
       \le\frac{2\pi |b|}{3M}
       \le\frac{8\,2^k}{M}.
 \tag{P-F56ah}
\]
The sum of $n-1-k$ over these bulk pairs is
$\binom{n-2}3-(n-4)$: before the one excluded pair,
it counts triples $2\le j<k<\ell\le n-1$.
There are $n-2$ exceptional odd factors, namely the
singleton one and the pairs using the zero-weight
variable. Each is at most $M/2$.
Indeed an odd residue modulo $M$ is at distance at least
one from $M/2$, so its cosine in (P-F56ag) has absolute
value at least $\sin(\pi/M)\ge2/M$, while its sine has
absolute value at most one.
Combining these estimates gives
\[
 \begin{split}
 \log_2\delta_n
 &\le 3\left(\binom{n-3}2-1\right)
        -\binom{n-2}3+(n-4)+(n-2)^2\\
 &=-n^3/6+4n^2-107n/6+19 .
 \end{split}
 \tag{P-F56ai}
\]
For $n\ge48$, $4n^2\le n^3/12$ and
$-107n/6+19<0$, proving the last assertion of
(P-F56ae).

Finally group the odd residues modulo $M_n$ into their
two lifts above an odd residue modulo $M_{n-1}$.
Reduction of $b_n$ is $b_{n-1}$.
The two selected parent classes therefore have the two
selected children $\pm b_n$ and the two unselected
children $\pm b_n+M_{n-1}$.
Complex conjugation makes the latter two contributions
equal, with total normalized mass $2\delta_n$.
For every other parent, (P-F56ac) and (P-F56ad)
bound its two children by $\rho$ times its normalized
mass in $Q_{n-1}$.
Here
$P_n(b_n)=J_n(b_n)P_{n-1}(b_{n-1})$,
using conjugation for the minus sign in (P-F56ac).
This proves (P-F56af), with no factor of two lost
in the nonselected branches.

Since $\delta_n\to0$ and $\rho<1$, the recursion forces
$Q_n\to0$. Explicitly, for any $\epsilon>0$, choose
$N$ so that $2\delta_n\le\epsilon(1-\rho)/2$
for all $n>N$. Iterating gives
\[
 Q_n\le \rho^{n-N}Q_N+\epsilon/2
                  \quad(n\ge N),
 \tag{P-F56aj}
\]
and the first term tends to zero.
The starting value $Q_N$ is finite by its definition.
This completes the conditional contraction reduction
and the unconditional selected-sibling bound.
